\section{Matrix extraction and Fourier decoding}
\label{sec:matrix}

Section~\ref{subsec:shortcode-proof} proves Theorem~\ref{thm:shortcode}.
The following subsections prove Proposition~\ref{prop:matrix-decoding},
with the parameters, table family, and advice law of
Section~\ref{sec:matrix-gap}, and Section~\ref{subsec:effective-k} closes
with the search for the tuple length.
Lemma~\ref{lem:matrix-spectral} first converts acceptance under the noise law $\mu$ into acceptance under uniform rank-one noise.  The inverse theorem
then supplies slices, and Lemma~\ref{lem:row-advice-agreement} makes slices with the required agreement
occur on a positive fraction of the visible row advice.
Lemma~\ref{lem:projected-witness} transfers a chosen slice to the projected
question.  Finally, Lemma~\ref{lem:fourier-coset-decoding} converts that
agreement into the probability of agreement of the two private decoders.

\subsection{Proof of the shortcode soundness theorem}\label{subsec:shortcode-proof}

We recall what Theorem~\ref{thm:shortcode} asserts.  Given
$\eta\in(0,1)$, we must find $\alpha\in(0,1]$, an integer
$r_{\rm s}\ge1$ and an integer $\ell_0$, depending only on $\eta$, and
for each $\ell\ge\ell_0$ a threshold $m_0(\eta,\ell)$, such that the
following holds for every $m\ge m_0$.  Every function
$f:\operatorname{Mat}_{\ell,m}\to\F_2^\ell$ whose equality acceptance
under the uniform rank-one step is at least $\eta$, as in
Equation~\eqref{eq:shortcode-acceptance} with $M\in\operatorname{Mat}_{\ell,m}$, $a\in\F_2^\ell$ and
$l\in\F_2^m$ independent and uniform, agrees with some affine
evaluation $Mz+u$ on at least an $\alpha$ fraction of some nonempty row
and column slice~\eqref{eq:shortcode-slice} with
$n_r,n_c\le r_{\rm s}$, as in Equation~\eqref{eq:shortcode-conclusion}.
No folding condition is assumed, and either factor of the step may be
zero.  The proof ends by making these choices effective, with rational
$\alpha$ for rational $\eta$, as asserted after the theorem.

\begin{proof}[Proof of Theorem~\ref{thm:shortcode}]
We derive this consequence from the Grassmann expansion theorem of
Khot--Minzer--Safra \cite[Theorem~1.12 in the 2018 revision]{KMS23},
using the matrix chart of Barak--Kothari--Steurer
\cite[Definition~3.2 and Lemmas~3.3--3.7]{BKS18}.

The Grassmann graph on an ambient binary space $E$ has the
$\ell$-dimensional subspaces as vertices; two vertices are adjacent when
their intersection has dimension $\ell-1$.  For a
nonempty vertex set, consider the probability that a uniform vertex in the set
has a uniform neighbor still in the set.  The cited theorem says that a
set with this probability at least a fixed $\zeta>0$ has positive relative
density in some interval
\[
 \{L:A_0\subseteq L\subseteq B_0\},
 \qquad \dim A_0+\operatorname{codim}_E B_0\le r.
\]
The density bound and $r$ depend only on $\zeta$.  The theorem holds for
all sufficiently large $\ell$, and then all sufficiently large
$\dim E$.

Equality acceptance is the average of the probabilities of remaining in the
nonempty fibers $f^{-1}(y)$ under the matrix step, weighted by their
uniform masses.  Some fiber $S$ therefore has probability at least $\eta$ of remaining in $S$ after one step.
The probability that either perturbation factor is zero is
\[
 p_0=2^{-\ell}+2^{-m}-2^{-\ell-m}.
\]
Choose the dimension thresholds so that $p_0\le\eta/2$.  Conditional on
both factors being nonzero, the probability of remaining in $S$ is at least
$(\eta-p_0)/(1-p_0)\ge\eta/2$.

For each matrix put
\[
 L_M=\{(x,M^\top x):x\in\F_2^\ell\}
 \subseteq E:=\F_2^\ell\oplus\F_2^m.
\]
The map $M\mapsto L_M$ is a bijection onto the chart $\mathcal C$ of
subspaces whose projection onto the first summand is invertible.  Since
\[
 \dim(L_M\cap L_N)=\ell-\rank(M-N),
\]
the neighbors within the chart correspond exactly to nonzero rank-one
matrix differences.  Each chart vertex has
$(2^\ell-1)(2^m-1)$ such neighbors, whereas its full Grassmann degree is
$(2^\ell-1)(2^{m+1}-2)$: choose a hyperplane in the vertex and then a
one-dimensional extension in the quotient by that hyperplane, other than
the vertex itself.  Thus exactly
half of its neighbors lie in the chart.  Over $\F_2$, every nonzero
rank-one matrix has a unique factorization $a l^\top$ with nonzero
factors.  The lifted set $\widetilde S=\{L_M:M\in S\}$ consequently has
probability at least $\eta/4$ of remaining in $\widetilde S$ after one Grassmann step.

Apply the expansion theorem with $\zeta=\eta/4$, obtaining density
$\alpha$ on an interval with
$\dim A_0+\operatorname{codim}_E B_0\le r_{\rm s}$.  Intersecting
this interval with $\mathcal C$ only increases the relative density of
$\widetilde S$, since $\widetilde S\subseteq\mathcal C$, and the
intersection is nonempty.  A basis $(d_i,s_i)$ of $A_0$ specifies
$A_0\subseteq L_M$ by the equations $d_i^\top M=s_i^\top$.
For the standard dot product on $E$, a basis $(t_j,q_j)$ of
$B_0^\perp$ specifies $L_M\subseteq B_0$ by
$Mq_j=t_j$.  Their simultaneous solution set is a slice of the required
form, with at most $r_{\rm s}$ equations of either kind.  On at least an
$\alpha$ fraction of this slice, $f$ equals the value defining $S$.
This proves the conclusion with $z=0$ and $u$ equal to that value.

The constants $\alpha,r_{\rm s}$ were chosen using only $\eta$.
Since $\dim E=\ell+m$, after fixing $\ell$ the remaining dimension
threshold is absorbed into $m_0(\eta,\ell)$.  No folding condition was
used.  The fiber argument is the equality consequence in
\cite[Lemma~2.3]{BKS18}; numbered references to that paper use its
arXiv version~1.

For the effective assertion, put $\zeta=\eta/4$.  The quantitative
argument in \cite[Section~2.4, Theorem~2.13, and Lemmas~2.7--2.8]{KMS23}
gives the following bound.  If a nonempty set $Y$ of vertices of the
Grassmann graph in ambient dimension $n$ has relative density at most
$\alpha$ on every interval
with $\dim A_0+\operatorname{codim}B_0\le r\le\ell$, then its
probability of remaining in the set is at most
\begin{equation}
 (r+1)2^{7r^3+3}(\alpha+2^{\ell+r-n})^{1/4}
       +2^{-(r+1)}+2^{-\ell}+2^{\ell-n+1}.
 \label{eq:effective-grassmann-staying-probability}
\end{equation}
To obtain this expression, lift $Y$ to the set $Y^*$ of all ordered bases
of its members, a basis-invariant set in the source's graph $H_{n,\ell}$.
This graph has vertex space $(\F_2^n)^\ell$ and step
$(x_i)_{i=1}^\ell\mapsto(x_i+b_i y)_{i=1}^\ell$, with independent
uniform $b\in\F_2^\ell$ and $y\in\F_2^n$.
The hypothesis says that $Y$ is $(r,\alpha)$-pseudorandom in the sense of
\cite[Definition~2.3]{KMS23}, so Lemma~2.8 there makes $Y^*$
$(r,\alpha+2^{\ell+r-n})$-pseudorandom in the sense of
\cite[Definition~2.5]{KMS23}.  This means that fixing any $q$ coordinates
and sampling the remaining coordinates independently and uniformly from
any subspace $W_0$ with $q+\operatorname{codim}W_0\le r$ gives membership
probability at most $\alpha+2^{\ell+r-n}$.

Let $F$ be the indicator of $Y^*$, let
$\delta>0$ be its density, and let $F_{=i}$ be the part of the Fourier
expansion of $F$ on the characters whose $\ell$ frequency vectors span a
space of dimension $i$; these characters have eigenvalue $2^{-i}$ under
the step \cite[Lemma~2.10]{KMS23}.  Theorem~2.13 there gives
$\norm{F_{=i}}_2^2\le2^{7r^3+3}(\alpha+2^{\ell+r-n})^{1/4}\delta$ for
$0\le i\le r$.  Parseval gives
$\sum_{i=0}^\ell\norm{F_{=i}}_2^2=\delta$.  Hence, as in
\cite[Section~2.4]{KMS23}, the probability $\rho^*$ of remaining in $Y^*$ after one step in $H_{n,\ell}$
satisfies
\begin{align*}
 \delta\rho^*
 &=\sum_i2^{-i}\norm{F_{=i}}_2^2
 \le\sum_{i\le r}\norm{F_{=i}}_2^2
   +2^{-(r+1)}\sum_{i>r}\norm{F_{=i}}_2^2\\
 &\le(r+1)2^{7r^3+3}(\alpha+2^{\ell+r-n})^{1/4}\delta
   +2^{-(r+1)}\delta.
\end{align*}
Divide by $\delta$.  Lemma~2.7 there bounds the difference between the
probabilities of remaining in $Y^*$ and $Y$ after one step by $2^{-\ell}+2^{\ell-n+1}$; adding it returns
to the Grassmann graph and gives
Equation~\eqref{eq:effective-grassmann-staying-probability}.

Choose $r\ge1$ with $2^{-(r+1)}\le\zeta/4$, and then a positive
rational $\alpha\le1$ with
\[
 2\alpha\bigl((r+1)2^{7r^3+3}\bigr)^4\le(\zeta/4)^4.
\]
Take $\ell\ge\max\{2,r\}$ with $2^{-\ell}\le\zeta/8$ and, after
fixing $\ell$, take $m$ so large that
$2^{r-m}\le\alpha$ and $2^{1-m}\le\zeta/8$.
With $n=\ell+m$, Equation~\eqref{eq:effective-grassmann-staying-probability}
is at most $3\zeta/4<\zeta$.  Thus a set with probability at least
$\zeta$ of remaining in the set after one step must have density greater than $\alpha$ on one of the
required intervals.  These same bounds make $p_0\le\eta/2$.
They give effective choices $r_{\rm s}=r$, $\ell_0$, and
$m_0(\eta,\ell)$ for the preceding argument, with the required
order of dependence.
\end{proof}

The next two subsections work for any positive integer $k$, with a fixed
family satisfying Equation~\eqref{eq:matrix-family-identities}.  They use
unrestricted uniform tables on $U$.
We specialize $k$ and $\beta$ when comparing that law with the projected
experiment.

\subsection{Comparing the two matrix perturbations}

The uniform matrix perturbation has the rank-indexed Fourier spectrum
computed in \cite[Lemma~2.10 in the 2018 revision]{KMS23}. We compare
that spectrum with the noise law $\mu$.

\begin{lemma}[Spectral comparison]
\label{lem:matrix-spectral}
If the table family passes Equation~\eqref{eq:matrix-test} with probability at
least $0.99$, then
\begin{equation}
 \E_{U,T}\Pr_{M,a,l}
   [F_U(M,T)=F_U(M+a l^\top,T)]
 \ge 0.9\cdot2^{-r_*}.
 \label{eq:shortcode-equality-acceptance}
\end{equation}
Here $U\sim\omega^{\otimes k}$; conditional on $U$, the maps
$T\in\Hom(L_U,W)$ and $M\in\Hom(L_U,K)$ and the vectors $a\in K$,
$l\in\F_2^m$ are independent and uniform.
\end{lemma}

\begin{proof}
For $\sigma\in K^*$ define the sign function
$f_{U,\sigma}(P)=(-1)^{\sigma(F_U(P))}$.  A frequency on the matrix
domain is a tuple $g=(g_i)_{i=0}^{m-1}$ of functionals in $\mathcal V^*$;
its character at $P$ is $(-1)^{\sum_i g_i(Pe_i)}$, where
$e_0,\ldots,e_{m-1}$ are the coordinate vectors of $L_U$.  All Fourier
coefficients use uniform probability measure.  Translation by $a l^\top$
multiplies this character by $(-1)^{\sum_i l_i g_i(a)}$.  Averaging over
uniform $l$ therefore gives the eigenvalue
\[
 \Lambda_g=\Pr_{a\sim\mu}[g_i(a)=0\text{ for every }i].
\]
It lies in $[0,1]$.  By Lemma~\ref{lem:equivariant-noise}, it is at most $7/8$
whenever
\[
 \dim\Span\{g_i|_K:0\le i<m\}\ge r_*.
\]
Indeed the span of the $g_i$ is a subspace of $\mathcal V^*$ to which
the lemma applies.

Expanding equality by a uniform output character, and then using Fourier
orthogonality, expresses the test acceptance as
\[
 \E_{U,\sigma}\sum_g
      \abs{\widehat f_{U,\sigma}(g)}^2\Lambda_g.
\]
Parseval gives total squared mass one for each sign function.  If $B$ is
the average squared mass on restricted ranks at least $r_*$, then
$0.99\le1-B/8$, so $B\le0.08$.  In particular, at least $0.9$ of the
average mass lies on ranks below $r_*$.

When $a$ is instead uniform in $K$, the eigenvalue at $g$ is
$2^{-\dim\Span\{g_i|_K\}}$.  The preceding mass thus gives the lower
bound in Equation~\eqref{eq:shortcode-equality-acceptance}.  This perturbation
leaves $T$ fixed, so the expression there is precisely the average of the
ordinary shortcode equality tests at fixed $(U,T)$.
\end{proof}

\subsection{Agreement on a positive fraction of row fibers}

The inverse theorem gives one structured slice on each qualifying
$(U,T)$ fiber.  The $U$-prover's information about the $K$-valued matrix
is the pair $(A,AM)$.  We need a positive probability that these observed rows determine
a slice with the required agreement.  We will also identify the short set of valid answers
supplied by that slice.  The erasure argument below proves the probability
bound, and folding then makes the answer set nonempty.

For this extraction step, sample $U$ as before and take independent uniform
maps $T:L_U\to W$, $A:K\to\F_2^{r_{\rm s}}$, and $M:L_U\to K$.
For fixed $U,T,A$ and a map
$S_0:L_U\to\F_2^{r_{\rm s}}$, the set $\{M:AM=S_0\}$ is its
\emph{row fiber}.  Let $\mathcal R$ consist of the row advice $(U,A,T,S_0)$ for which there exist
at most $r_{\rm s}$ column specifications and a target $Mz+u$ such that
the simultaneous solution set of the column specifications and $AM=S_0$
is nonempty and $F_U(M,T)=Mz+u$ with probability at least $\alpha/2$
under the uniform measure on this set.

This is an analytic condition on the row fiber.  The target vector $z$
need not initially be an actual answer.  The fixed column values allow
$z$ to vary modulo the span of their column directions, with a corresponding
change in $u$.  Thus the slice naturally supplies a coset of possible
target vectors.  Its valid part will be the answer set sampled by the
$U$-prover.

\begin{lemma}[Row advice admitting agreement and its short answer set]
\label{lem:row-advice-agreement}
Fix the parameters of Equation~\eqref{eq:matrix-parameters} and a table
family satisfying Equation~\eqref{eq:matrix-family-identities}.
For every sufficiently large positive integer $k$, acceptance at least
$0.99$ in the matrix test implies
\begin{equation}
 \Pr_{U,T,A,M}[(U,A,T,AM)\in\mathcal R]
 \ge g_0,
 \qquad
 g_0=10\eta^2\,2^{-\ell r_{\rm s}}>0.
 \label{eq:advice-agreement-mass}
\end{equation}
The lower bound and the required largeness of $k$ are independent of the
parity instance and the family.

For every positive integer $k$, each row advice in $\mathcal R$ admits a
deterministically chosen certificate consisting of at most $r_{\rm s}$
column specifications
$Mq_j=t_j$, with $q_j\in L_U$ and $t_j\in K$, and a target $Mz+u$,
with $z\in L_U$, $u\in K$, and $e_0^\top z=1$, satisfying the definition
of $\mathcal R$.
Writing $Z=\Span\{q_j\}$, the chosen witness determines the answer set
\begin{equation}
 \mathcal A=(z+Z)\cap\{v\in L_U:e_0^\top v=1\},
 \qquad 1\le |\mathcal A|\le 2^{r_{\rm s}}.
 \label{eq:row-answer-set}
\end{equation}
The witness and $\mathcal A$ depend only on $(U,A,T,S_0)$ and the fixed
table family.  The slice and its target retain the agreement guarantee in
the definition of $\mathcal R$.
\end{lemma}

\begin{proof}
By Lemma~\ref{lem:matrix-spectral}, the average shortcode equality acceptance
is at least $0.9\cdot2^{-r_*}=14.4\eta$.  If $q$ is the probability that
the acceptance at $(U,T)$ is at least $4\eta$, the average is at most
$4\eta+q$, so $q\ge10.4\eta\ge10\eta$.
Fix such a pair $(U,T)$ and write $f(M)=F_U(M,T)$.
Let $B$ be the union, in the matrix domain, of all
row fibers of advice in $\mathcal R$, over every possible $A,S_0$.  For each $M\in B$, at
least one of the $2^{\ell r_{\rm s}}$ maps $A$ gives advice in $\mathcal R$ on the row fiber
containing $M$, namely the fiber of $S_0=AM$.  The uniform choice of
$A$ hits that map with probability $2^{-\ell r_{\rm s}}$, so at this
$(U,T)$ the probability over $A,M$ that $(U,A,T,AM)\in\mathcal R$ is at least
$2^{-\ell r_{\rm s}}\Pr[B]$.  It therefore suffices to show that
$\Pr[B]\ge\eta$.
We will erase all row fibers of advice in $\mathcal R$.  If their union $B$ were too small,
acceptance would survive the erasure, while independent random
replacement values would destroy every possible agreement certificate.
Suppose that its uniform measure is less than $\eta$.  Change $f$ at the entries of $B$, replacing
each by an independent uniform element of $K$.  Each matrix receives one
replacement, even when several such row fibers overlap there.  Both endpoints of the
shortcode step are uniform matrices.  Consequently every such modification
changes acceptance by at most $2\Pr[B]$, and the modified function has
acceptance at least $2\eta$.

Put $r=r_{\rm s}$.  We claim that the replacement values can be chosen
so that the following holds simultaneously for every nonempty slice
$\mathcal S$ with at most $r$ row and $r$ column specifications and
every target $Mz+u$, with $z\in\F_2^m$ and $u\in K$: fewer than
$(\alpha/4)|\mathcal S|$ matrices $M\in\mathcal S\cap B$ receive
$Mz+u$ as their replacement value.  The space of differences of matrices in such a slice is
\[
 \Hom\bigl(\F_2^m/\Span\{q_j\},\ \bigcap_i\ker d_i^\top\bigr).
\]
Its size $n_{\mathcal S}$ is therefore at least
$2^{(\ell-r)(m-r)}$.  Padding the lists of specifications with zero
equations, the number of slice--target descriptions is at most
$2^{(2r+1)(\ell+m)}$.

For a fixed description, count only matches at randomized entries and
assign a deterministic zero contribution to the other entries.  These are
independent variables in $[0,1]$, with total mean at most
$2^{-\ell}n_{\mathcal S}<\alpha n_{\mathcal S}/8$.  Hoeffding's
inequality \cite[Theorem~2, equation~(2.6)]{Hoeffding1963} bounds the
probability of at least
$\alpha n_{\mathcal S}/4$ such matches by
$\exp(-\alpha^2 n_{\mathcal S}/32)$: apply it to the average of
$n_{\mathcal S}$ summands with deviation at least $\alpha/8$,
giving exponent $-2n_{\mathcal S}(\alpha/8)^2$.
The union bound over all
descriptions is thus at most
\[
 \exp\left((2r+1)(\ell+m)\log2
       -\frac{\alpha^2}{32}2^{(\ell-r)(m-r)}\right),
\]
which tends to zero as $m=1+2k$ grows.  Fix replacements with the asserted
property.

Consider any slice of the allowed type.  If its row advice was in $\mathcal R$, every
entry of that fiber belongs to $B$, and its modified agreement is less
than $\alpha/4$.  If the row advice was outside $\mathcal R$, the original agreement
with each target on the slice was less than $\alpha/2$.  The unmodified
matches form a subset of those original matches, and the randomized
matches add less than $\alpha/4$.  The modified agreement is then less
than $3\alpha/4$.  Padding shorter row systems by zeros shows that these
two cases cover every slice in Theorem~\ref{thm:shortcode}.  That theorem
now contradicts the modified function's acceptance at least $2\eta$.
Its applicability does not require the modification to preserve folding.

It follows that $\Pr[B]\ge\eta$, so at this fixed $(U,T)$ the
probability that $(U,A,T,AM)\in\mathcal R$ is at least
$\eta2^{-\ell r_{\rm s}}$.  At least $10\eta$ of the pairs qualify,
as shown at the start of the proof, proving
Equation~\eqref{eq:advice-agreement-mass}.  All counting and concentration
bounds depended only on the fixed parameters and $m$, as required.

The remaining assertion does not use high test acceptance or largeness of
$k$.  Fix any positive integer $k$, a family satisfying
Equation~\eqref{eq:matrix-family-identities}, and any row advice in $\mathcal R$.
Fix a witness, write its column specifications as $Mq_j=t_j$, and put
$Z=\Span\{q_j\}$.  Nonemptiness of the slice makes the prescribed values
consistent: they determine a linear map $t:Z\to K$ with $t(w)=Mw$
throughout the slice.  Replacing $z$ by $z+w$ and $u$ by $u+t(w)$ for
$w\in Z$ therefore preserves the target.  It suffices to show that the
coset $z+Z$ contains a vector of first bit one.

Otherwise $z$ and every vector in $Z$ have first bit zero.  Put $H=\ker A$.
For each $h\in H$, translation $M\mapsto M+h e_0^\top$ preserves the
row fiber, the column specifications, and the target.  By folding it
translates $F_U(M,T)$ by $h$.  The action is free, and in each of its
$H$-orbits at most one matrix can satisfy the target equality.  The
agreement is at most
\[
 |H|^{-1}=2^{-\ell+\rank A}
 \le2^{-(\ell-r_{\rm s})}<\alpha/8,
\]
contradicting membership in $\mathcal R$.  Choose $w\in Z$ for which $z+w$ has first bit one
and make the corresponding changes to $z$ and $u$.  They preserve the
original agreement fraction.  Choose a witness and such a representative
using a fixed order on the finite possibilities.  This choice uses only
the row advice and the fixed table family.  Its vector $z$ lies in
$\mathcal A$, so the set is nonempty; since $\dim Z\le r_{\rm s}$,
its size is at most $2^{r_{\rm s}}$.
\end{proof}

\subsection{Transferring the chosen slice}

The comparison with sparsely projected questions follows the smooth
outer-game method of Khot--Safra~\cite[Section~3 and Lemma~3.1]{KS13},
with shared linear advice as in Khot--Minzer--Safra~\cite[Sections~3.3--3.4]{KMS17}.
Lemma~\ref{lem:row-advice-agreement} gives a positive probability of row advice
that determines a slice and a short set of valid answers under unrestricted
sampling on $U$.  In the actual experiment, the maps are sampled on $O$
and padded to $U$.  We compare the complete padded table with the
unrestricted table, including the affine intercept, with projection choices
averaged out, and then transfer the chosen slice to $O$.

From now on take $k$ along integer cubes larger than one and set
$\beta=k^{-2/3}$ in Definition~\ref{def:projected-advice}.

The choice $\beta=k^{-2/3}$ serves two purposes: $k\beta^2\to0$ makes
the padded table law approach the unrestricted law in
Lemma~\ref{lem:projected-tv}, while $k\beta\to\infty$ makes the expected
number of singleton coordinates diverge.  Once the advice spaces are
fixed, each singleton has a uniformly positive probability that every
visible advice row has zero coefficient on its variable bit.
Lemma~\ref{lem:zero-advice} uses these coordinates to bound agreement.

\begin{lemma}[The padded matrix is asymptotically uniform]
\label{lem:projected-tv}
For each equation question $U$, the padded table
\[
 (M,T)=(M_O\bar\pi,T_O\bar\pi)
\]
in Definition~\ref{def:projected-advice}, with $\beta=k^{-2/3}$, has
conditional distribution at total variation distance $o(1)$ from the unrestricted uniform
distribution on $\Hom(L_U,K)\times\Hom(L_U,W)$, as $k$ tends to
infinity along cubes.  The bound is uniform in $U$ and the parity instance.
The projection choices are averaged out in this comparison.
\end{lemma}

\begin{proof}
Combine the two components into a $\mathcal V$-valued table and put
$n=|\mathcal V|$.  On an unprojected equation block the two slopes are
independent uniform elements of $\mathcal V$.  On a singleton block the
padded pair of slopes is an equal mixture of
\[
 (a,0),\qquad (0,a),\qquad (a,a),
 \qquad a\text{ uniform in }\mathcal V.
\]
Call these pair laws $P_0$ and $Q_0$, respectively.  The likelihood ratio
$dQ_0/dP_0$ is $n$ at $(0,0)$, $n/3$ at the other $3(n-1)$ points on
the three displayed lines, and zero elsewhere.  Thus
\[
 \chi^2(Q_0\Vert P_0)
 :=\E_{P_0}\left[\left(\frac{dQ_0}{dP_0}-1\right)^2\right]
 =\frac{n-1}{3}.
\]
The one-block mixture $(1-\beta)P_0+\beta Q_0$ consequently has
chi-squared divergence $\beta^2(n-1)/3$ from $P_0$.  Blocks are sampled
independently, including when their question IDs coincide.  Multiplying
their likelihood ratios gives divergence
\begin{equation}
 \left(1+\frac{n-1}{3}\beta^2\right)^k-1=o(1),
 \label{eq:projected-chi-square}
\end{equation}
because $n$ is fixed and $k\beta^2=k^{-1/3}$ tends to zero.  Total
variation is at most one half the square root of this expression.

For a projection onto the third bit of equation $e$, the padded intercept
acquires the term $b_e a$, where $a$ is that singleton's slope.  The
original intercept is independent and uniform in $\mathcal V$.  Adding
the sum of these terms leaves it uniform and independent even after all
projection choices and slopes are given.  The preceding slope estimate
therefore proves the assertion for the entire table.  Neither the slope
laws nor the bound in Equation~\eqref{eq:projected-chi-square} depend on
the equation IDs or their right-hand sides.
\end{proof}

Fix $U$-advice in $\mathcal R$ for a moment, with its chosen witness and answer set
$\mathcal A$ from Lemma~\ref{lem:row-advice-agreement}.  For a coordinate
projection put $z'=\bar\pi z$ and $Z'=\bar\pi Z$.  Since $\bar\pi$
preserves the first coordinate,
\begin{equation}
 \bar\pi\mathcal A
 =(z'+Z')\cap\{v\in L_O:e_0^\top v=1\}.
 \label{eq:projected-answer-set}
\end{equation}
Indeed, every vector in the coset on the right has a preimage in $z+Z$,
and the preimage has the same first coordinate.  We will show that the
$O$-prover's private Fourier sample reaches this projected answer set.

The certifying slice does not automatically survive projection.  For
example, a column condition $Mq=t$ becomes $M_O\bar\pi q=t$; it is
inconsistent if $\bar\pi q=0$ but $t\ne0$.  More generally, projected
column directions may acquire new dependencies.  The next argument uses
the comparison of the full table laws to find a positive probability of samples
where the chosen slice remains nonempty and retains enough agreement.

For this analysis, write
$\mathsf D=(U,O,\pi,A,T_O,S_O)$ for the joint data that determine both prover
messages and the projection.  These data do not specify $M_O$ beyond the
constraint $AM_O=S_O$.

\begin{lemma}[Transfer of a chosen witness]
\label{lem:projected-witness}
Fix a table family satisfying Equation~\eqref{eq:matrix-family-identities}
and having matrix-test acceptance at least $0.99$.  In the projected
advice experiment with $\beta=k^{-2/3}$, for every sufficiently large
integer cube $k$, there is a set of values of the joint data $\mathsf D$, of
total probability at least
\begin{equation}
 \gamma_1=\frac{\alpha g_0}{8}\,2^{-\ell r_{\rm s}}>0
 \label{eq:projected-witness-mass}
\end{equation}
such that every $\mathsf D$ in this set has the following properties.
The padded $U$-advice belongs to $\mathcal R$, with the chosen answer set $\mathcal A$ and its
certifying witness.  Writing $z'=\bar\pi z$, that witness transfers to a
nonempty slice
\[
 AM_O=S_O,\qquad M_O\bar\pi q_j=t_j,
\]
on which
\begin{equation}
 \Pr\bigl[F_O(M_O,T_O)=M_Oz'+u\mid
           AM_O=S_O,\ M_O\bar\pi q_j=t_j\text{ for all }j\bigr]
 \ge\alpha/4.
 \label{eq:transferred-witness}
\end{equation}
The probability in this equation is taken over a uniform matrix in the
specified slice.
\end{lemma}

\begin{proof}
First use unrestricted sampling of $U,A,M,T$, setting $S_0=AM$.  Let
$\mathcal J$ be the event that $(U,A,T,S_0)\in\mathcal R$ and $M$ satisfies
its chosen column specifications.  Let $\mathcal M$ be equality with
the chosen target, declaring this event false for row advice outside $\mathcal R$.  Within a
row fiber of advice in $\mathcal R$ the nonempty column slice
imposes at most $\ell r_{\rm s}$ additional scalar linear equations.
Hence Lemma~\ref{lem:row-advice-agreement} gives
\[
 \Pr[\mathcal J]\ge g_0\,2^{-\ell r_{\rm s}},
 \qquad
 \Pr[\mathcal J\cap\mathcal M]\ge(\alpha/2)\Pr[\mathcal J].
\]
It follows that
\begin{equation}
 \Pr[\mathcal J\cap\mathcal M]
       -(\alpha/4)\Pr[\mathcal J]
 \ge\frac{\alpha g_0}{4}\,2^{-\ell r_{\rm s}}.
 \label{eq:witness-positive-difference}
\end{equation}

Both events in this expression depend only on $(U,A,M,T)$.  Adding the
independent $A$ to Lemma~\ref{lem:projected-tv} preserves its variation
bound.  We may therefore compare the two event probabilities with their
values in the padded experiment.  If the total variation distance is
$v_k=o(1)$, the difference in Equation~\eqref{eq:witness-positive-difference}
changes by at most $(1+\alpha/4)v_k$.  For all sufficiently large $k$
the padded difference is at least $\gamma_1$.

Now condition on $\mathsf D$.  The remaining $M_O$ is uniform on the row fiber
$AM_O=S_O$.  The padded
$U$-advice, and therefore its chosen witness and answer set, are fixed.
On data with $U$-advice in $\mathcal R$ let $p_{\mathsf D}$ be the conditional probability of the projected
column slice, and, when $p_{\mathsf D}>0$, let $q_{\mathsf D}$ be its conditional equality
fraction.  Put $p_{\mathsf D}=0$ on data with $U$-advice outside $\mathcal R$, and put $q_{\mathsf D}=0$ when
$p_{\mathsf D}=0$.  Sharing gives
$F_U(M_O\bar\pi,T_O\bar\pi)=F_O(M_O,T_O)$, so the padded difference
is exactly
\[
 \E_{\mathsf D}\bigl[p_{\mathsf D}(q_{\mathsf D}-\alpha/4)\bigr]\ge\gamma_1.
\]
The integrand is nonpositive outside the event consisting of $U$-advice in $\mathcal R$,
$p_{\mathsf D}>0$, and $q_{\mathsf D}\ge\alpha/4$, and it is at most one everywhere.
That event thus has probability at least $\gamma_1$ and has precisely the
asserted properties.  The law comparison applied to events depending only
on $(U,A,M,T)$; conditioning on $\mathsf D$ here only identifies where their
positive difference occurs in the projected experiment.
\end{proof}

\subsection{Decoding the row advice}

Extracting an outer strategy from Fourier mass follows the approach of
H\aa stad~\cite[Section~1.1]{Hastad01} and the Fourier-list method of
Khot--Safra~\cite[Section~1.1]{KS13}. Here the projected prover samples
squared Fourier mass on its own row fiber.  The transferred slice will
show that this private sample reaches the projected answer set
$\bar\pi\mathcal A$ with positive probability.

We now define both strategies from their separate messages.  The joint
data and the transferred witness serve only in the analysis of their
agreement; the $O$-prover does not use the $U$-prover's witness or answer set.

The $U$-prover uses its visible row advice to make the deterministic choice
in Lemma~\ref{lem:row-advice-agreement}.  When the advice belongs to $\mathcal R$, it samples
$a_U$ uniformly from $\mathcal A$.
This is a nonempty set of at most $2^{r_{\rm s}}$ actual answers.  On
other inputs the prover returns any actual answer.

The $O$-prover starts from its own row fiber.  Given $(O,A,T_O,S_O)$, put
$H=\ker A$ and choose, by a fixed rule, a representative $M_0$ satisfying
$AM_0=S_O$.  Such a representative exists
on every input in the experiment.  Write the row fiber as $M_0+\mathcal N$, where
$\mathcal N=\Hom(L_O,H)$.  Sample uniform $\sigma\in K^*$ and form the sign
function
\[
 f_\sigma(N)=(-1)^{\sigma(F_O(M_0+N,T_O))},\qquad N\in \mathcal N.
\]
Sample a Fourier frequency $\Phi\in \mathcal N^*$ with probability
$|\widehat f_\sigma(\Phi)|^2$.  Identify
$\mathcal N^*=H^*\otimes L_O$, with the pairing defined on pure tensors by
\[
 \langle\varphi\otimes v,N\rangle=\varphi(Nv).
\]
For $h\in H$, contraction sends $\Phi$ to a vector $\Phi(h)\in L_O$.
If $\sigma|_H\ne0$, choose by a fixed rule an $h\in H$ with $\sigma(h)=1$ and
return $\Phi(h)$.  If the restriction is trivial, return any actual
answer.

This rule always returns an actual answer at frequencies of positive
mass.  Indeed the folding identity in
Equation~\eqref{eq:matrix-family-identities} gives
\[
 f_\sigma(N+h e_0^\top)=-f_\sigma(N)
 \qquad\text{when }\sigma(h)=1.
\]
Changing variables in a Fourier coefficient shows that its being nonzero
requires
$(-1)^{\langle\Phi,h e_0^\top\rangle}=-1$, equivalently
$e_0^\top\Phi(h)=1$.

Figure~\ref{fig:advice} summarizes the information and private choices of
the two decoders.
\begin{figure}[htbp]
\centering
\begin{tikzpicture}[
  x=1cm,y=1cm,
  every node/.style={font=\small},
  box/.style={draw=linkblue!80!black,rounded corners=2pt,
    text width=6.0cm,align=center,inner sep=8pt},
  input/.style={box,minimum height=16mm},
  rule/.style={box,minimum height=35mm},
  output/.style={box,minimum height=10mm},
  flow/.style={-{Latex[length=2.1mm]},thick,draw=black!70}
]
\node[box,text width=13.3cm,minimum height=12mm] (sample) at (3.7,0)
  {\textbf{Referee}\\[2pt]
   Sample $U,O,\bar\pi,A,T_O,M_O$; set $S_O=AM_O$.
   Send only the inputs below.};

\node[input] (uinput) at (0,-1.9)
  {\textbf{$U$-decoder} (equation question)\\
   $(U,A,T_U,S_U)$\\
   $T_U=T_O\bar\pi,\quad S_U=S_O\bar\pi$};
\node[input] (oinput) at (7.4,-1.9)
  {\textbf{$O$-decoder} (projected question)\\
   $(O,A,T_O,S_O)$};
\draw[flow] (sample.south -| uinput.north) -- (uinput.north);
\draw[flow] (sample.south -| oinput.north) -- (oinput.north);

\node[rule] (urule) at (0,-4.8)
  {Choose the fixed slice witness and\\its set of valid answers\\[2pt]
   $\mathcal A=(z+Z)\cap\{v:e_0^\top v=1\}$\\
   from this decoder's visible advice.\\[6pt]
   \textit{Private randomness:}\\
   sample $a_U$ uniformly from $\mathcal A$.};
\node[rule] (orule) at (7.4,-4.8)
  {Set $H=\ker A$. Choose $M_0$\\with $AM_0=S_O$.\\[5pt]
   \textit{Private randomness:}\\
   sample uniform $\sigma\in K^*$, then $\Phi$\\
   with probability $|\widehat f_\sigma(\Phi)|^2$.\\[5pt]
   If $\sigma|_H\ne0$, use a fixed $h\in H$\\
   with $\sigma(h)=1$; set $a_O=\Phi(h)$.};
\draw[flow] (uinput) -- (urule);
\draw[flow] (oinput) -- (orule);

\node[output] (uoutput) at (0,-7.9)
  {$a_U\in L_U,\qquad e_0^\top a_U=1$};
\node[output] (ooutput) at (7.4,-7.9)
  {$a_O\in L_O,\qquad e_0^\top a_O=1$};
\draw[flow] (urule) -- (uoutput);
\draw[flow] (orule) -- (ooutput);

\node[box,text width=8.2cm,minimum height=11mm] (verify) at (3.7,-9.6)
  {\textbf{Referee checks}\quad $\bar\pi a_U=a_O$};
\draw[flow] (uoutput.south) -- (verify.north -| uoutput.south);
\draw[flow] (ooutput.south) -- (verify.north -| ooutput.south);
\end{tikzpicture}
\caption{The two decoders use independent random tapes. Neither receives the
sampled $M_O$ or the projection $\bar\pi$; the $O$-decoder chooses its own
representative $M_0$ from the visible advice. Every displayed advice map
includes its affine intercept, through the coordinate $t_0$. The
transferred slice is used only to analyze these private rules: the
$O$-decoder does not use the $U$-decoder's witness or answer set. If the
$U$-advice is outside $\mathcal R$, or if $\sigma|_H=0$ for the $O$-decoder, that
decoder returns an arbitrary valid answer.}
\label{fig:advice}
\end{figure}

\FloatBarrier

\begin{lemma}[A transferred slice gives agreement]
\label{lem:fourier-coset-decoding}
Fix joint data $\mathsf D$ for which the chosen witness has the nonempty
projected slice and agreement in Equation~\eqref{eq:transferred-witness}.
Conditional on $\mathsf D$, the two decoders defined above, with independent
randomness, satisfy $\bar\pi a_U=a_O$ with probability at least
\begin{equation}
 2^{-\ell r_{\rm s}}(\alpha/8)^2\,2^{-r_{\rm s}}.
 \label{eq:conditional-decoding}
\end{equation}
\end{lemma}

\begin{proof}
Put $Z'=\bar\pi Z$, so $\dim Z'\le r_{\rm s}$.  The projected column
specifications in the row fiber fix the restriction $N|_{Z'}$ to a linear
map $b:Z'\to H$.  The resulting affine slice is nonempty; choose one
$N_*$ in it.  For every $\sigma\in K^*$ define
\[
 C_\sigma=\E\bigl[f_\sigma(N)(-1)^{\sigma(Nz')}
                       \mid N|_{Z'}=b\bigr].
\]
The character expansion of Equation~\eqref{eq:transferred-witness} gives
\[
 \E_\sigma\bigl[(-1)^{\sigma(M_0z'+u)}C_\sigma\bigr]
 \ge\alpha/4,
 \qquad\text{so}\qquad
 \E_\sigma|C_\sigma|\ge\alpha/4.
\]
Exactly a fraction $2^{-\dim H}$ of the characters have trivial
restriction to $H$.  Since $\dim H\ge\ell-r_{\rm s}$ and
$|C_\sigma|\le1$, Equation~\eqref{eq:matrix-parameters} implies
\begin{equation}
 \E_\sigma\bigl[\mathbf 1_{\{\sigma|_H\ne0\}}|C_\sigma|\bigr]
 \ge\alpha/8.
 \label{eq:nontrivial-correlation}
\end{equation}

The annihilator in $\mathcal N^*$ of matrices vanishing on $Z'$ is $H^*\otimes Z'$.
Indeed, restriction $\mathcal N\to\Hom(Z',H)$ is onto, since a map on $Z'$
extends after choosing a basis of $L_O$.  Its dual image is precisely
the functionals depending only on this restriction, generated by
$N\mapsto\varphi(Nv)$ for $\varphi\in H^*$ and $v\in Z'$.
Writing $Y=\{N\in \mathcal N:N|_{Z'}=0\}$, the slice is $N_*+Y$.
For any frequency $\Xi\in \mathcal N^*$, orthogonality on $Y$ gives
\[
 \E_{N\in N_*+Y}(-1)^{\langle\Xi,N\rangle}
 =\mathbf1_{\{\Xi\in H^*\otimes Z'\}}
       (-1)^{\langle\Xi,N_*\rangle}.
\]
For $\theta_\sigma=(\sigma|_H)\otimes z'$, Fourier expansion on the
affine slice therefore gives
\[
 C_\sigma
 =\sum_{\Phi\in\theta_\sigma+H^*\otimes Z'}
       \widehat f_\sigma(\Phi)
       (-1)^{\langle\Phi+\theta_\sigma,N_*\rangle}.
\]
This sum has $2^{\dim H\dim Z'}\le2^{\ell r_{\rm s}}$ terms.
Cauchy--Schwarz, first on each sum and then in the average over $\sigma$,
shows from Equation~\eqref{eq:nontrivial-correlation} that
\begin{align*}
 &\E_\sigma\left[
   \mathbf 1_{\{\sigma|_H\ne0\}}
   \sum_{\Phi\in\theta_\sigma+H^*\otimes Z'}
             |\widehat f_\sigma(\Phi)|^2\right]\\
 &\hspace{15mm}\ge
 2^{-\ell r_{\rm s}}
 \E_\sigma\bigl[\mathbf 1_{\{\sigma|_H\ne0\}}C_\sigma^2\bigr]
 \ge 2^{-\ell r_{\rm s}}(\alpha/8)^2.
\end{align*}
This is a lower bound on the probability that the $O$-decoder samples a
nontrivial restriction and a frequency in the indicated coset.

Every such frequency has the form
$\Phi=(\sigma|_H)\otimes z'+\Psi$ with $\Psi\in H^*\otimes Z'$.
For the decoder's choice of $h$ it satisfies
\[
 \Phi(h)=z'+\Psi(h)\in z'+Z'.
\]
As already proved, a frequency of positive mass gives first bit one.
Equation~\eqref{eq:projected-answer-set} therefore places the returned
vector in $\bar\pi\mathcal A$.  Independently, the $U$-decoder samples
uniformly from $\mathcal A$ and chooses a preimage with probability at
least $1/|\mathcal A|\ge2^{-r_{\rm s}}$.  Multiplying the bounds proves
Equation~\eqref{eq:conditional-decoding}.
\end{proof}

\begin{proof}[Proof of Proposition~\ref{prop:matrix-decoding}]
Use the two decoders above.  Lemma~\ref{lem:projected-witness} supplies a
set of joint data of probability at least $\gamma_1$, and
Lemma~\ref{lem:fourier-coset-decoding} gives the stated conditional
agreement on that set.  Thus one may take
\[
 \gamma=\gamma_1\,2^{-\ell r_{\rm s}}(\alpha/8)^2\,2^{-r_{\rm s}}>0.
\]
The largeness requirements for $k$ come only from
Theorem~\ref{thm:shortcode}, the concentration bound in
Lemma~\ref{lem:row-advice-agreement}, and the total variation estimate in
Lemma~\ref{lem:projected-tv}.  They are uniform in the parity instance
and table family.
\end{proof}

\subsection{Choosing the tuple length}\label{subsec:effective-k}

The threshold in Proposition~\ref{prop:matrix-gap} can be met by a
terminating search over cubes once the map, noise, and shortcode data are fixed.
The constants $g_0$, $\gamma_1$, and $\gamma$ have been specified in
Equation~\eqref{eq:advice-agreement-mass},
Equation~\eqref{eq:projected-witness-mass}, and the proof just given.
Write $m=1+2k$ and $\beta=k^{-2/3}$; since $k$ is a cube, $\beta$ is
rational.  Require $m\ge m_0(\eta,\ell)$ and
\[
 \frac{\alpha^2}{32}\,2^{(\ell-r_{\rm s})(m-r_{\rm s})}
 >(2r_{\rm s}+1)(\ell+m).
\]
The latter is a sufficient rational condition for the union bound in
Lemma~\ref{lem:row-advice-agreement} to be less than one.  For witness
transfer in Lemma~\ref{lem:projected-witness}, require
\[
 \frac14\left[
 \left(1+\frac{|\mathcal V|-1}{3}\beta^2\right)^k-1\right]
 <\frac{\gamma_1^2}{(1+\alpha/4)^2},
\]
where $\gamma_1$ is from
Equation~\eqref{eq:projected-witness-mass}.  The left-hand side is the
square of the total-variation bound.  For the final contradiction,
require
\[
 \left(1-\frac{\beta\,2^{-\dim W-r_{\rm s}}}{3600}\right)^k
 <\gamma,
\]
with $\gamma$ from Proposition~\ref{prop:matrix-decoding}.  All these
conditions are rational and hold for sufficiently large cubes $k$.
Enumerating cubes until they hold therefore finds a suitable $k$;
the search cost depends only on the two prescribed errors and is
incurred before the input formula is used.

